\subsection{FOURIER TRANSFORMS OF CENTRAL FUNCTIONS}

For all \(u \in \hat{\mathrm{G}}\) , denote by \(\chi_u\) the character of \(u\) ; thus,

\[
\chi_ {u} (g) = \operatorname{Tr} (u (g)), \quad (g \in \mathrm{G}).\tag{16}
\]

Recall from Spectral Theory the formulas

\[
\chi_ {u} * \chi_ {v} = 0 (u, v \in \hat {\mathrm{G}}, u \neq v)\tag{17}
\]

\[
\chi_ {u} * \chi_ {u} = \frac {1}{d (u)} \chi_ {u} \quad (u \in \hat {\mathrm{G}}).\tag{18}
\]

For all \(u \in \hat{G}\) , denote by \(\varepsilon_{u}\) the identity map of \(E_{u}\) . Recall (§7, no. 4) that \(\mathrm{ZL}^{2}(\mathrm{G})\) denotes the subspace of \(\mathrm{L}^{2}(\mathrm{G})\) consisting of the classes of the functions f that are central, that is, such that \(f \circ \operatorname{Int}s = f\) for all \(s \in G\) , or equivalently that \(\gamma(s)f = \delta(s^{-1})f\) for all \(s \in G\) .

PROPOSITION 4. Let \(f \in \mathrm{L}^2(\mathrm{G})\) . Then \(f\) is central if and only if \(u(f)\) is a homothety for all \(u \in \hat{\mathrm{G}}\) . In that case

\[
u (f) = \frac {\varepsilon_ {u}}{d (u)} \int_ {\mathrm{G}} f (g) \chi_ {u} (g) d g.\tag{19}
\]

By Prop. 1 (no. 1), to say that \(f\) is central means that \(u(\gamma(s)f) = u(\delta(s^{-1})f)\) for all \(s \in \mathbf{G}\) and all \(u \in \hat{\mathbf{G}}\) ; but this can also be written as \(u(s)u(f) = u(f)u(s)\) for all \(s \in \mathbf{G}\) and all \(u \in \hat{\mathbf{G}}\) (formulas (12) and (13)), hence the first assertion of Prop. 4 (Schur's lemma). If \(u(f)\) is a homothety, then \(u(f) = \lambda_u\varepsilon_u\) with

\[
\lambda_ {u} = \frac {1}{d (u)} \operatorname{Tr} (u (f)) = \frac {1}{d (u)} \int_ {\mathrm{G}} f (g) \operatorname{Tr} (u (g)) d g = \frac {1}{d (u)} \int_ {\mathrm{G}} f (g) \chi_ {u} (g) d g.
\]

Consequently, for \(f \in \mathrm{ZL}^2(\mathbf{G})\) we have

\[
u (f) = \langle \overline {{\chi}} _ {u} | f \rangle \frac {\varepsilon_ {u}}{d (u)} \rangle ,\tag{20}
\]

SO

\[
\overline {{{{\mathcal {F}}}}} (f) = \left(\langle \overline {{{{\chi}}}} _ {u} | f \rangle \frac {\varepsilon_ {u}}{d (u)}\right) _ {u \in \hat {\mathrm{G}}}\tag{21}
\]

with

\[
\| \overline {{\mathcal {F}}} (f) \| _ {2} ^ {2} = \sum_ {u} \left| \left| \langle \overline {{\chi}} _ {u} | f \rangle \frac {\varepsilon_ {u}}{d (u)} \right| \right| _ {2} ^ {2} = \sum_ {u} | \langle \overline {{\chi}} _ {u} | f \rangle | ^ {2}.
\]

Conversely, if \(\varphi\) is a square-integrable complex function on \(\hat{\mathbf{G}}\) , the element \(\left(\frac{\varphi(u)}{d(u)}\varepsilon_u\right)_{u\in \hat{\mathbf{G}}}\) of \(\mathrm{F}(\hat{\mathrm{G}})\) belongs to \(\mathrm{L}^2 (\hat{\mathrm{G}})\) , and we have (formula (9))

\[
\left(\mathcal {F} _ {u} \left(\frac {\varphi (u)}{d (u)} \varepsilon_ {u}\right)\right) (g) = d (u) \mathrm{Tr} \left(\frac {\varphi (u)}{d (u)} \varepsilon_ {u} u (g) ^ {- 1}\right) = \varphi (u) \overline {{\chi}} _ {u} (g),
\]

SO

\[
\mathcal {F} \left(\left(\frac {\varphi (u)}{d (u)} \varepsilon_ {u}\right)\right) = \sum_ {u \in \hat {\mathrm{G}}} \varphi (u) \overline {{\chi}} _ {u}.\tag{22}
\]

Note, in particular, that formulas (20) and (21) give, for u, v in \(\hat{G}\) ,

\[
u (\overline {{\chi}} _ {v}) = 0 \quad \text { if } u \neq v,\tag{23}
\]

\[
u (\overline {{\chi}} _ {u}) = \frac {\varepsilon_ {u}}{d (u)} \in \mathrm{End} (\mathrm{E} _ {u}),\tag{24}
\]

\[
\overline {{\mathcal {F}}} (\chi_ {u}) = \frac {\varepsilon_ {u}}{d (u)} \in \mathrm{End} (\mathrm{E} _ {u}) \subset \mathrm{F} (\hat {\mathrm{G}}).\tag{25}
\]

PROPOSITION 5. Let f be a continuous central function on G. Then f is infinitely-differentiable if and only if the function \(u \mapsto |\langle \chi_{u}|f\rangle|\) is rapidly decreasing on \(\hat{G}\) ; in that case,

\[
f (g) = \sum_ {u \in \hat {\mathrm{G}}} \langle \chi_ {u} | f \rangle \chi_ {u} (g)
\]

for all g ∈ G.

By Th. 1 b), the function \(\overline{f}\) is infinitely-differentiable if and only if the function \(u \mapsto \| u(\overline{f})\|_{\infty}\) is rapidly decreasing; but, by (20),

\[
\left\| u (\overline {{f}}) \right\| _ {\infty} = \frac {1}{d (u)} | \langle \chi_ {u} | f \rangle |,
\]

hence the first assertion, since the functions \(d(u)\) and \(\frac{1}{d(u)}\) are moderately increasing.

Assume that \(f\) is infinitely-differentiable; by Th. 1 a), \(f(g) = \sum_{u\in \hat{\mathbf{G}}}f_u(g)\) for all \(g\in \mathrm{G}\) , so

\[
\begin{array}{c} f _ {u} (g) = \langle u (g) | u (f) \rangle = d (u) \mathrm{Tr} (u (g) ^ {- 1}. u (f)) = d (u) \mathrm{Tr} \left(u (g) ^ {- 1} \langle \overline {{\chi}} _ {u} | f \rangle \frac {\varepsilon_ {u}}{d (u)}\right) \\ = \langle \overline {{\chi}} _ {u} | f \rangle \mathrm{Tr} (u (g) ^ {- 1}) = \langle \overline {{\chi}} _ {u} | f \rangle \overline {{\chi}} _ {u} (g). \end{array}
\]

Hence, \(f(g) = \sum_{u\in \hat{\mathbf{G}}}\langle \overline{\chi}_u|f\rangle \overline{\chi}_u(g)\) ; but, for all \(u\in \hat{\mathbf{G}}\) , the contragredient representation \(u'\) of \(u\) satisfies \(\overline{\chi}_u = \chi_{u'}\) and the map \(u\mapsto u'\) is a permutation of \(\hat{\mathbf{G}}\) ; so we also have \(f(g) = \sum_{u\in \hat{\mathbf{G}}}\langle \chi_u|f\rangle \chi_u(g)\) , hence the proposition.

COROLLARY. Let \(f\) be a continuous central function on \(G\) . Then \(f\) is infinitely-differentiable if and only if the restriction of \(f\) to \(T\) is infinitely-differentiable.

Indeed, by Cor. 4 of §7, no. 4,

\[
\langle \chi_ {u} | f \rangle = \int_ {\mathrm{G}} \overline {{\lambda (u) (t)}} \varphi (t) d t, \quad \text { where } \varphi (t) = \prod_ {\alpha > 0} (1 - \alpha (t) ^ {- 1}) f (t).
\]

If \(f|_{\mathrm{T}}\) is infinitely-differentiable, so is \(\varphi\) ; by Prop. 5, applied to the group \(\mathrm{T}\) , the function \(\mu \mapsto \int_{\mathrm{T}} \overline{\mu(t)} \varphi(t) dt\) on \(\hat{\mathrm{T}} = \mathrm{X}(\mathrm{T})\) is then rapidly decreasing, and so is the function \(u \mapsto \langle \chi_u | f \rangle\) ; hence the function \(f\) is infinitely-differentiable (Prop. 5). The converse is clear.

